% Options for packages loaded elsewhere
\PassOptionsToPackage{unicode}{hyperref}
\PassOptionsToPackage{hyphens}{url}
\documentclass[
]{article}
\usepackage{xcolor}
\usepackage[margin=1in]{geometry}
\usepackage{amsmath,amssymb}
\setcounter{secnumdepth}{-\maxdimen} % remove section numbering
\usepackage{iftex}
\ifPDFTeX
  \usepackage[T1]{fontenc}
  \usepackage[utf8]{inputenc}
  \usepackage{textcomp} % provide euro and other symbols
\else % if luatex or xetex
  \usepackage{unicode-math} % this also loads fontspec
  \defaultfontfeatures{Scale=MatchLowercase}
  \defaultfontfeatures[\rmfamily]{Ligatures=TeX,Scale=1}
\fi
\usepackage{lmodern}
\ifPDFTeX\else
  % xetex/luatex font selection
\fi
% Use upquote if available, for straight quotes in verbatim environments
\IfFileExists{upquote.sty}{\usepackage{upquote}}{}
\IfFileExists{microtype.sty}{% use microtype if available
  \usepackage[]{microtype}
  \UseMicrotypeSet[protrusion]{basicmath} % disable protrusion for tt fonts
}{}
\makeatletter
\@ifundefined{KOMAClassName}{% if non-KOMA class
  \IfFileExists{parskip.sty}{%
    \usepackage{parskip}
  }{% else
    \setlength{\parindent}{0pt}
    \setlength{\parskip}{6pt plus 2pt minus 1pt}}
}{% if KOMA class
  \KOMAoptions{parskip=half}}
\makeatother
\usepackage{color}
\usepackage{fancyvrb}
\newcommand{\VerbBar}{|}
\newcommand{\VERB}{\Verb[commandchars=\\\{\}]}
\DefineVerbatimEnvironment{Highlighting}{Verbatim}{commandchars=\\\{\}}
% Add ',fontsize=\small' for more characters per line
\usepackage{framed}
\definecolor{shadecolor}{RGB}{248,248,248}
\newenvironment{Shaded}{\begin{snugshade}}{\end{snugshade}}
\newcommand{\AlertTok}[1]{\textcolor[rgb]{0.94,0.16,0.16}{#1}}
\newcommand{\AnnotationTok}[1]{\textcolor[rgb]{0.56,0.35,0.01}{\textbf{\textit{#1}}}}
\newcommand{\AttributeTok}[1]{\textcolor[rgb]{0.13,0.29,0.53}{#1}}
\newcommand{\BaseNTok}[1]{\textcolor[rgb]{0.00,0.00,0.81}{#1}}
\newcommand{\BuiltInTok}[1]{#1}
\newcommand{\CharTok}[1]{\textcolor[rgb]{0.31,0.60,0.02}{#1}}
\newcommand{\CommentTok}[1]{\textcolor[rgb]{0.56,0.35,0.01}{\textit{#1}}}
\newcommand{\CommentVarTok}[1]{\textcolor[rgb]{0.56,0.35,0.01}{\textbf{\textit{#1}}}}
\newcommand{\ConstantTok}[1]{\textcolor[rgb]{0.56,0.35,0.01}{#1}}
\newcommand{\ControlFlowTok}[1]{\textcolor[rgb]{0.13,0.29,0.53}{\textbf{#1}}}
\newcommand{\DataTypeTok}[1]{\textcolor[rgb]{0.13,0.29,0.53}{#1}}
\newcommand{\DecValTok}[1]{\textcolor[rgb]{0.00,0.00,0.81}{#1}}
\newcommand{\DocumentationTok}[1]{\textcolor[rgb]{0.56,0.35,0.01}{\textbf{\textit{#1}}}}
\newcommand{\ErrorTok}[1]{\textcolor[rgb]{0.64,0.00,0.00}{\textbf{#1}}}
\newcommand{\ExtensionTok}[1]{#1}
\newcommand{\FloatTok}[1]{\textcolor[rgb]{0.00,0.00,0.81}{#1}}
\newcommand{\FunctionTok}[1]{\textcolor[rgb]{0.13,0.29,0.53}{\textbf{#1}}}
\newcommand{\ImportTok}[1]{#1}
\newcommand{\InformationTok}[1]{\textcolor[rgb]{0.56,0.35,0.01}{\textbf{\textit{#1}}}}
\newcommand{\KeywordTok}[1]{\textcolor[rgb]{0.13,0.29,0.53}{\textbf{#1}}}
\newcommand{\NormalTok}[1]{#1}
\newcommand{\OperatorTok}[1]{\textcolor[rgb]{0.81,0.36,0.00}{\textbf{#1}}}
\newcommand{\OtherTok}[1]{\textcolor[rgb]{0.56,0.35,0.01}{#1}}
\newcommand{\PreprocessorTok}[1]{\textcolor[rgb]{0.56,0.35,0.01}{\textit{#1}}}
\newcommand{\RegionMarkerTok}[1]{#1}
\newcommand{\SpecialCharTok}[1]{\textcolor[rgb]{0.81,0.36,0.00}{\textbf{#1}}}
\newcommand{\SpecialStringTok}[1]{\textcolor[rgb]{0.31,0.60,0.02}{#1}}
\newcommand{\StringTok}[1]{\textcolor[rgb]{0.31,0.60,0.02}{#1}}
\newcommand{\VariableTok}[1]{\textcolor[rgb]{0.00,0.00,0.00}{#1}}
\newcommand{\VerbatimStringTok}[1]{\textcolor[rgb]{0.31,0.60,0.02}{#1}}
\newcommand{\WarningTok}[1]{\textcolor[rgb]{0.56,0.35,0.01}{\textbf{\textit{#1}}}}
\usepackage{graphicx}
\makeatletter
\newsavebox\pandoc@box
\newcommand*\pandocbounded[1]{% scales image to fit in text height/width
  \sbox\pandoc@box{#1}%
  \Gscale@div\@tempa{\textheight}{\dimexpr\ht\pandoc@box+\dp\pandoc@box\relax}%
  \Gscale@div\@tempb{\linewidth}{\wd\pandoc@box}%
  \ifdim\@tempb\p@<\@tempa\p@\let\@tempa\@tempb\fi% select the smaller of both
  \ifdim\@tempa\p@<\p@\scalebox{\@tempa}{\usebox\pandoc@box}%
  \else\usebox{\pandoc@box}%
  \fi%
}
% Set default figure placement to htbp
\def\fps@figure{htbp}
\makeatother
\setlength{\emergencystretch}{3em} % prevent overfull lines
\providecommand{\tightlist}{%
  \setlength{\itemsep}{0pt}\setlength{\parskip}{0pt}}
\usepackage{bookmark}
\IfFileExists{xurl.sty}{\usepackage{xurl}}{} % add URL line breaks if available
\urlstyle{same}
\hypersetup{
  pdftitle={Bayesian Mediation Analysis},
  hidelinks,
  pdfcreator={LaTeX via pandoc}}

\title{Bayesian Mediation Analysis}
\author{}
\date{\vspace{-2.5em}2026-04-17}

\begin{document}
\maketitle

\subsection{Introduction}\label{introduction}

Mediation analysis decomposes a treatment effect into the part that
flows through an intermediate variable (the mediator) and the part that
does not. Frequentist approaches build confidence intervals around the
product \(a_1 b_2\) via Sobel tests or non-parametric bootstrapping;
both have to compensate for the fact that the product of two normally
distributed estimates is itself non-normal. A Bayesian formulation
sidesteps the issue: the joint posterior of \((a_1, b_2)\) is sampled
directly, the indirect effect is computed draw-by-draw as the product,
and any quantile of the resulting posterior is a valid credible interval
without distributional assumptions.

\subsection{Prerequisites}\label{prerequisites}

A working understanding of multiple regression, the concept of mediation
in causal inference, and Bayesian regression. Familiarity with
multivariate \texttt{brms} formulas accelerates the implementation step.

\subsection{Theory}\label{theory}

The classical Baron-Kenny system has two regressions

\[M = a_0 + a_1 X + \varepsilon_M, \qquad Y = b_0 + b_1 X + b_2 M + \varepsilon_Y,\]

with errors that are typically taken Normal. The direct effect of \(X\)
on \(Y\) is \(b_1\), the indirect effect through \(M\) is \(a_1 b_2\),
and the total effect is \(b_1 + a_1 b_2\). In the Bayesian formulation,
the two equations are fitted jointly so that the posterior carries the
joint uncertainty in \((a_1, b_2)\). The indirect-effect posterior is
the empirical distribution of \(\{a_1^{(s)} b_2^{(s)}\}_{s=1}^S\); the
credible interval comes from its quantiles. No symmetry assumption is
needed and the posterior captures any skew induced by the product.

\subsection{Assumptions}\label{assumptions}

Linearity of both equations, Normal residuals (or whatever family you
specify), no unmeasured confounders of the \(M \to Y\) relation, and
correctly ordered causal direction \(X \to M \to Y\). Mediation analysis
is a causal claim; the Bayesian machinery does not make weak design
choices stronger.

\subsection{R Implementation}\label{r-implementation}

\begin{Shaded}
\begin{Highlighting}[]
\FunctionTok{library}\NormalTok{(brms)}

\FunctionTok{set.seed}\NormalTok{(}\DecValTok{2026}\NormalTok{)}
\NormalTok{n }\OtherTok{\textless{}{-}} \DecValTok{300}
\NormalTok{X }\OtherTok{\textless{}{-}} \FunctionTok{rnorm}\NormalTok{(n)}
\NormalTok{M }\OtherTok{\textless{}{-}} \FloatTok{0.5} \SpecialCharTok{*}\NormalTok{ X }\SpecialCharTok{+} \FunctionTok{rnorm}\NormalTok{(n)}
\NormalTok{Y }\OtherTok{\textless{}{-}} \FloatTok{0.3} \SpecialCharTok{*}\NormalTok{ X }\SpecialCharTok{+} \FloatTok{0.6} \SpecialCharTok{*}\NormalTok{ M }\SpecialCharTok{+} \FunctionTok{rnorm}\NormalTok{(n)}
\NormalTok{d }\OtherTok{\textless{}{-}} \FunctionTok{data.frame}\NormalTok{(X, M, Y)}

\CommentTok{\# Joint model using brms multivariate syntax}
\NormalTok{bf1 }\OtherTok{\textless{}{-}} \FunctionTok{bf}\NormalTok{(M }\SpecialCharTok{\textasciitilde{}}\NormalTok{ X)}
\NormalTok{bf2 }\OtherTok{\textless{}{-}} \FunctionTok{bf}\NormalTok{(Y }\SpecialCharTok{\textasciitilde{}}\NormalTok{ X }\SpecialCharTok{+}\NormalTok{ M)}

\NormalTok{fit }\OtherTok{\textless{}{-}} \FunctionTok{brm}\NormalTok{(bf1 }\SpecialCharTok{+}\NormalTok{ bf2 }\SpecialCharTok{+} \FunctionTok{set\_rescor}\NormalTok{(}\ConstantTok{FALSE}\NormalTok{), }\AttributeTok{data =}\NormalTok{ d,}
           \AttributeTok{chains =} \DecValTok{4}\NormalTok{, }\AttributeTok{iter =} \DecValTok{2000}\NormalTok{, }\AttributeTok{refresh =} \DecValTok{0}\NormalTok{)}

\NormalTok{post }\OtherTok{\textless{}{-}} \FunctionTok{as\_draws\_df}\NormalTok{(fit)}
\NormalTok{indirect }\OtherTok{\textless{}{-}}\NormalTok{ post}\SpecialCharTok{$}\NormalTok{b\_M\_X }\SpecialCharTok{*}\NormalTok{ post}\SpecialCharTok{$}\NormalTok{b\_Y\_M}
\FunctionTok{quantile}\NormalTok{(indirect, }\FunctionTok{c}\NormalTok{(}\FloatTok{0.025}\NormalTok{, }\FloatTok{0.5}\NormalTok{, }\FloatTok{0.975}\NormalTok{))}
\end{Highlighting}
\end{Shaded}

\subsection{Output \& Results}\label{output-results}

A multivariate \texttt{brms} fit jointly estimates the two regressions
and stores draws for both. Computing the elementwise product of
posterior draws of \(a_1\) and \(b_2\) produces the indirect-effect
posterior; quantiles of that posterior give the credible interval
directly. The total effect, the direct effect, and the proportion
mediated all come from the same draw matrix without re-fitting.

\subsection{Interpretation}\label{interpretation}

A reporting sentence: ``Bayesian mediation estimated the indirect effect
of \(X\) on \(Y\) through \(M\) at 0.30 (95 \% credible interval 0.18 to
0.45) and the direct effect at 0.28 (95 \% CrI 0.14 to 0.42); the
indirect path explained approximately 52 \% of the total effect (95 \%
CrI 35 \% to 68 \%).'' Always report both direct and indirect effects
with credible intervals; reporting only the indirect product hides
whether the direct path remains substantive.

\subsection{Practical Tips}\label{practical-tips}

\begin{itemize}
\tightlist
\item
  Use multivariate \texttt{brms} formulas
  (\texttt{bf(M\ \textasciitilde{}\ X)\ +\ bf(Y\ \textasciitilde{}\ X\ +\ M)\ +\ set\_rescor(FALSE)})
  so that the two regressions are fitted jointly and the joint posterior
  is preserved; running them separately throws away the cross-equation
  correlation needed for honest indirect-effect intervals.
\item
  Place priors on each regression's coefficients individually; weakly
  informative defaults (\(\mathcal N(0, 2.5)\) on standardised scales)
  are usually appropriate.
\item
  Sensitivity analysis to unmeasured confounding is more important than
  the choice of estimator; consider \texttt{causalsens} or Bayesian
  instrumental-variable extensions when the assumption is weak.
\item
  For binary or count mediators / outcomes, change the family in the
  corresponding \texttt{bf()} and the indirect-effect computation
  generalises; just use \texttt{posterior\_epred()} to translate to the
  natural-effects scale.
\item
  Reporting the proportion mediated (\(a_1 b_2 / (b_1 + a_1 b_2)\))
  requires care because the denominator can be small or change sign
  across draws; truncate or report the indirect and total effects
  separately when the proportion is unstable.
\end{itemize}

\subsection{R Packages Used}\label{r-packages-used}

\texttt{brms} for joint multivariate regressions and direct posterior
products of coefficients, \texttt{bmlm} for purpose-built Bayesian
mediation with multilevel structure, \texttt{tidybayes} for tidy
posterior draws of the indirect effect, and \texttt{bayestestR} for
credible intervals and ROPE summaries on the derived quantities.

\subsection{For Reviewers}\label{for-reviewers}

\textbf{What to look for in a paper using this method.}

\begin{itemize}
\tightlist
\item
  \textbf{Common misapplications.}
\item
  \textbf{Diagnostics that should be reported but often aren't.}
\item
  \textbf{Red flags in tables and figures.}
\item
  \textbf{What to verify.}
\item
  \textbf{What an adequate Methods paragraph must contain.}
\end{itemize}

\subsection{See also --- labs in this
chapter}\label{see-also-labs-in-this-chapter}

\begin{itemize}
\tightlist
\item
  \href{labs/lab-bayesian-thinking-control-vs-decision-error.Rmd}{Bayesian
  thinking; α control vs decision error}
\item
  \href{labs/lab-brms-stan-loo-hierarchical-models.Rmd}{brms / Stan,
  LOO, hierarchical models}
\end{itemize}

\end{document}
